% Options for packages loaded elsewhere
\PassOptionsToPackage{unicode}{hyperref}
\PassOptionsToPackage{hyphens}{url}
\PassOptionsToPackage{dvipsnames,svgnames,x11names}{xcolor}
\documentclass[
  10pt,
]{article}
\usepackage{xcolor}
\usepackage[margin=1.8cm]{geometry}
\usepackage{amsmath,amssymb}
\setcounter{secnumdepth}{-\maxdimen} % remove section numbering
\usepackage{iftex}
\ifPDFTeX
  \usepackage[T1]{fontenc}
  \usepackage[utf8]{inputenc}
  \usepackage{textcomp} % provide euro and other symbols
\else % if luatex or xetex
  \usepackage{fontspec}
  \defaultfontfeatures{Scale=MatchLowercase}
  \defaultfontfeatures[\rmfamily]{Ligatures=TeX,Scale=1}
\fi
\usepackage{lmodern}
\ifPDFTeX\else
  % xetex/luatex font selection
\fi
% Use upquote if available, for straight quotes in verbatim environments
\IfFileExists{upquote.sty}{\usepackage{upquote}}{}
\IfFileExists{microtype.sty}{% use microtype if available
  \usepackage[]{microtype}
  \UseMicrotypeSet[protrusion]{basicmath} % disable protrusion for tt fonts
}{}
\makeatletter
\@ifundefined{KOMAClassName}{% if non-KOMA class
  \IfFileExists{parskip.sty}{%
    \usepackage{parskip}
  }{% else
    \setlength{\parindent}{0pt}
    \setlength{\parskip}{6pt plus 2pt minus 1pt}}
}{% if KOMA class
  \KOMAoptions{parskip=half}}
\makeatother
\usepackage{color}
\usepackage{fancyvrb}
\newcommand{\VerbBar}{|}
\newcommand{\VERB}{\Verb[commandchars=\\\{\}]}
\DefineVerbatimEnvironment{Highlighting}{Verbatim}{commandchars=\\\{\}}
% Add ',fontsize=\small' for more characters per line
\newenvironment{Shaded}{}{}
\newcommand{\AlertTok}[1]{\textcolor[rgb]{1.00,0.00,0.00}{\textbf{#1}}}
\newcommand{\AnnotationTok}[1]{\textcolor[rgb]{0.38,0.63,0.69}{\textbf{\textit{#1}}}}
\newcommand{\AttributeTok}[1]{\textcolor[rgb]{0.49,0.56,0.16}{#1}}
\newcommand{\BaseNTok}[1]{\textcolor[rgb]{0.25,0.63,0.44}{#1}}
\newcommand{\BuiltInTok}[1]{\textcolor[rgb]{0.00,0.50,0.00}{#1}}
\newcommand{\CharTok}[1]{\textcolor[rgb]{0.25,0.44,0.63}{#1}}
\newcommand{\CommentTok}[1]{\textcolor[rgb]{0.38,0.63,0.69}{\textit{#1}}}
\newcommand{\CommentVarTok}[1]{\textcolor[rgb]{0.38,0.63,0.69}{\textbf{\textit{#1}}}}
\newcommand{\ConstantTok}[1]{\textcolor[rgb]{0.53,0.00,0.00}{#1}}
\newcommand{\ControlFlowTok}[1]{\textcolor[rgb]{0.00,0.44,0.13}{\textbf{#1}}}
\newcommand{\DataTypeTok}[1]{\textcolor[rgb]{0.56,0.13,0.00}{#1}}
\newcommand{\DecValTok}[1]{\textcolor[rgb]{0.25,0.63,0.44}{#1}}
\newcommand{\DocumentationTok}[1]{\textcolor[rgb]{0.73,0.13,0.13}{\textit{#1}}}
\newcommand{\ErrorTok}[1]{\textcolor[rgb]{1.00,0.00,0.00}{\textbf{#1}}}
\newcommand{\ExtensionTok}[1]{#1}
\newcommand{\FloatTok}[1]{\textcolor[rgb]{0.25,0.63,0.44}{#1}}
\newcommand{\FunctionTok}[1]{\textcolor[rgb]{0.02,0.16,0.49}{#1}}
\newcommand{\ImportTok}[1]{\textcolor[rgb]{0.00,0.50,0.00}{\textbf{#1}}}
\newcommand{\InformationTok}[1]{\textcolor[rgb]{0.38,0.63,0.69}{\textbf{\textit{#1}}}}
\newcommand{\KeywordTok}[1]{\textcolor[rgb]{0.00,0.44,0.13}{\textbf{#1}}}
\newcommand{\NormalTok}[1]{#1}
\newcommand{\OperatorTok}[1]{\textcolor[rgb]{0.40,0.40,0.40}{#1}}
\newcommand{\OtherTok}[1]{\textcolor[rgb]{0.00,0.44,0.13}{#1}}
\newcommand{\PreprocessorTok}[1]{\textcolor[rgb]{0.74,0.48,0.00}{#1}}
\newcommand{\RegionMarkerTok}[1]{#1}
\newcommand{\SpecialCharTok}[1]{\textcolor[rgb]{0.25,0.44,0.63}{#1}}
\newcommand{\SpecialStringTok}[1]{\textcolor[rgb]{0.73,0.40,0.53}{#1}}
\newcommand{\StringTok}[1]{\textcolor[rgb]{0.25,0.44,0.63}{#1}}
\newcommand{\VariableTok}[1]{\textcolor[rgb]{0.10,0.09,0.49}{#1}}
\newcommand{\VerbatimStringTok}[1]{\textcolor[rgb]{0.25,0.44,0.63}{#1}}
\newcommand{\WarningTok}[1]{\textcolor[rgb]{0.38,0.63,0.69}{\textbf{\textit{#1}}}}
\usepackage{longtable,booktabs,array}
\newcounter{none} % for unnumbered tables
\usepackage{calc} % for calculating minipage widths
% Correct order of tables after \paragraph or \subparagraph
\usepackage{etoolbox}
\makeatletter
\patchcmd\longtable{\par}{\if@noskipsec\mbox{}\fi\par}{}{}
\makeatother
% Allow footnotes in longtable head/foot
\IfFileExists{footnotehyper.sty}{\usepackage{footnotehyper}}{\usepackage{footnote}}
\makesavenoteenv{longtable}
\setlength{\emergencystretch}{3em} % prevent overfull lines
\providecommand{\tightlist}{%
  \setlength{\itemsep}{0pt}\setlength{\parskip}{0pt}}
\usepackage{bookmark}
\IfFileExists{xurl.sty}{\usepackage{xurl}}{} % add URL line breaks if available
\urlstyle{same}
\hypersetup{
  pdftitle={Response to JC's review: what he asked{,} what we found{,} what we change},
  colorlinks=true,
  linkcolor={Maroon},
  filecolor={Maroon},
  citecolor={Blue},
  urlcolor={Blue},
  pdfcreator={LaTeX via pandoc}}

\title{Response to JC's review: what he asked, what we found, what we
change}
\author{}
\date{}

\begin{document}
\maketitle

Draft, 19 Nov 2026. For Narmeen, as the source for the meeting slides
and for the paper edits.

\textbf{How to read the references.}

\begin{itemize}
\tightlist
\item
  \textbf{``p.~4, Sec. 4 `Selection status'\,''} = the PDF JC read
  (``TRDP\_Paper (2).pdf'', 9 pages, attached with his email).
\item
  \textbf{\texttt{sec/4\_protocol.tex:10}} = the same text in the repo
  copy of the paper (\texttt{sd-lora-attention-segmentation/paper/}),
  where the edit will be made. The repo copy matches JC's PDF except for
  the equation numbering in Sec. 3.3 (see \#6).
\item
  \textbf{{[}pending{]}} = a run is still on the cluster; the cell is
  filled in when it finishes (same day).
\end{itemize}

\textbf{Words used throughout.}

\begin{itemize}
\tightlist
\item
  \textbf{BRACE}: our full method (SA-guided CA seed \ensuremath{\rightarrow} coarse-to-fine SA
  with the one-sided cap \ensuremath{\rightarrow} DenseCRF \ensuremath{\rightarrow} threshold 0.55).
\item
  \textbf{Fine-only}: the simpler pipeline without the coarse-to-fine
  hierarchy (the seed is propagated twice through the fine 64-pixel SA
  only, then DenseCRF).
\item
  \textbf{mIoU}: how well masks overlap the annotation, 0--1, averaged
  over the 20 classes.
\item
  \textbf{VOC-Sim}: 490 simple images, one object (``a photograph of a
  cat''). \textbf{COCO-Cap}: 707 cluttered images from COCO captions.
\item
  \textbf{Out-of-fold (OOF)}: VOC-Sim is split into 5 parts; each choice
  is made on 4 parts and scored on the 5th. \textbf{Forced OOF}: one
  fixed variant scored that way.
\item
  \textbf{VOC-sel.}: the threshold is chosen on VOC-Sim and reused
  unchanged on COCO-Cap.
\item
  \textbf{Paired}: two methods scored on exactly the same images; the
  difference (\ensuremath{\Delta}) comes with a 95\% interval in brackets. If the interval
  does not contain 0, the difference is real.
\item
  \textbf{DenseCRF}: a standard post-processing step that snaps the mask
  to image edges.
\end{itemize}

\begin{center}\rule{0.5\linewidth}{0.5pt}\end{center}

\section{\#1 Contribution and claim}

\begin{quote}
\emph{JC:} ``In Table 1, the two new components (SA-guided layer
weighting and the capped hierarchy) add +0.006 each on VOC-Sim; most of
the gain over OVAM comes from configuration choices (CA scales, SA
passes, DenseCRF, call window). The strongest argument for the hierarchy
is COCO-Cap (+0.034), i.e.~cluttered scenes. I would frame the
contribution around that, rather than around \guillemotleft{}+12.5 points over OVAM\guillemotright{}.
Also note that both components were adopted after the selection rule had
kept the simpler reference (Sec. 4): this is honest, but reviewers will
read it against the contribution \guillemotleft{}fold-safe selection protocol\guillemotright{}, so the
text must handle it carefully.''
\end{quote}

\subsection{1.1 What he means, in plain words}

\textbf{Part 1: the headline is the wrong number.} The paper sells BRACE
by its margin over OVAM (+0.125 on VOC-Sim, +0.114 on COCO-Cap). Table 1
shows where that margin comes from:

{\def\LTcaptype{none} % do not increment counter
\begin{longtable}[]{@{}
  >{\raggedright\arraybackslash}p{(\linewidth - 6\tabcolsep) * \real{0.2500}}
  >{\raggedright\arraybackslash}p{(\linewidth - 6\tabcolsep) * \real{0.2500}}
  >{\raggedright\arraybackslash}p{(\linewidth - 6\tabcolsep) * \real{0.2500}}
  >{\raggedright\arraybackslash}p{(\linewidth - 6\tabcolsep) * \real{0.2500}}@{}}
\toprule\noalign{}
\begin{minipage}[b]{\linewidth}\raggedright
Table 1 row (p.~4)
\end{minipage} & \begin{minipage}[b]{\linewidth}\raggedright
VOC-Sim
\end{minipage} & \begin{minipage}[b]{\linewidth}\raggedright
\ensuremath{\Delta}
\end{minipage} & \begin{minipage}[b]{\linewidth}\raggedright
What kind of change
\end{minipage} \\
\midrule\noalign{}
\endhead
\bottomrule\noalign{}
\endlastfoot
CA only, all scales & .665 & & starting point \\
drop the 32-pixel CA layers & .691 & +.026 & configuration \\
add 64-pixel SA propagation & .789 & +.033\textsuperscript{a} (really .691 \ensuremath{\rightarrow} .789 = +.098)
& configuration \\
two passes & .800 & +.012 & configuration \\
DenseCRF & .813 & +.014 & configuration \\
SA calls 10--20 & .817 & +.004 & configuration \\
\textbf{SA-guided layer weights (new)} & .823 & \textbf{+.006} &
\textbf{our new component} \\
\textbf{coarse-to-fine + cap (new)} & .829 & \textbf{+.006} &
\textbf{our new component} \\
\end{longtable}
}

So on VOC-Sim, about 0.15 of the gain is configuration (choices anyone
could make on top of OVAM) and only 0.012 is our two new ideas. A
reviewer will say: ``the novelty adds 1 point''. \textbf{But on COCO-Cap
the new components matter:} .644 \ensuremath{\rightarrow} .662 (+0.018, layer weights) \ensuremath{\rightarrow} .696
(+0.034, hierarchy). JC's advice: make the paper about \emph{why the
scale structure matters in cluttered scenes}, not about beating OVAM.

\textbf{Part 2: the selection rule and the ``simpler reference''.} You
asked what these are.

\begin{itemize}
\tightlist
\item
  \textbf{The selection rule} is written in Sec. 4, ``Fold-safe
  selection'' (p.~4, \texttt{sec/4\_protocol.tex:7}): a change is
  accepted only if it gains at least 0.010 mIoU on the training folds
  (or 0.005 with at least 12 of 20 classes better and \textbf{no class
  losing more than 0.05}), and wins at least 4 of 5 folds.
\item
  \textbf{The simpler reference} is the pipeline \emph{without} the new
  component. There are two:
\end{itemize}

{\def\LTcaptype{none} % do not increment counter
\begin{longtable}[]{@{}
  >{\raggedright\arraybackslash}p{(\linewidth - 4\tabcolsep) * \real{0.3333}}
  >{\raggedright\arraybackslash}p{(\linewidth - 4\tabcolsep) * \real{0.3333}}
  >{\raggedright\arraybackslash}p{(\linewidth - 4\tabcolsep) * \real{0.3333}}@{}}
\toprule\noalign{}
\begin{minipage}[b]{\linewidth}\raggedright
New component (the ``complicated'' choice)
\end{minipage} & \begin{minipage}[b]{\linewidth}\raggedright
Simpler reference the rule kept
\end{minipage} & \begin{minipage}[b]{\linewidth}\raggedright
Why the rule kept the simpler one
\end{minipage} \\
\midrule\noalign{}
\endhead
\bottomrule\noalign{}
\endlastfoot
SA-guided CA layer weights & \textbf{uniform} layer weights (plain
average over CA layers) & Won 4/5 folds (+0.006), but the
multi-candidate screen kept uniform weights (supp. Tab. 10, p.~10 of the
supplement: ``selector kept uniform weights'') \\
coarse-to-fine hierarchy with the cap (BRACE) & \textbf{fine-only}
propagation + DenseCRF & Won 5/5 folds (+0.0063), but one class,
\textbf{sheep, lost 0.057--0.104}, more than the 0.05 guard \\
\end{longtable}
}

So the rule, applied as registered, would have given a paper whose final
method is \textbf{fine-only with uniform weights} (0.823 on VOC-Sim).
The owners overrode the rule and adopted both components. The paper says
so in one sentence (p.~4, Sec. 4 ``Selection status'',
\texttt{sec/4\_protocol.tex:10}) and once more in Sec. 7 (p.~8,
\texttt{sec/7\_discussion.tex:7}).

\subsection{1.2 The problem, exactly}

\begin{enumerate}
\def\labelenumi{\arabic{enumi}.}
\tightlist
\item
  \textbf{The paper contradicts itself.} Contribution 4 (p.~1,
  \texttt{sec/1\_intro.tex:14}) says ``every choice made on VOC-Sim
  folds'', and the intro says ``Every choice is made on VOC-Sim folds,
  with COCO-Cap used only after freezing'' (p.~1,
  \texttt{sec/1\_intro.tex:7}). But the two choices that \emph{are} the
  novelty were not made by the rule.
\item
  \textbf{The override saw COCO-Cap.} The protocol log (Section 135, 31
  Oct) shows the decision to adopt the hierarchy was taken
  \textbf{after} COCO-Cap was scored, and it cites the COCO-Cap +0.0341
  as evidence. So COCO-Cap is not an untouched test set for that
  decision, yet the paper presents COCO-Cap as ``frozen transfer'' and
  JC himself calls it the strongest argument. A reviewer who notices
  this will read it as choosing the method on the test set.
\item
  \textbf{``Out-of-fold'' is used for a number the procedure did not
  produce.} 0.829 is a \emph{forced} OOF score of the owner-chosen
  variant (abstract, intro). The out-of-fold score of the registered
  procedure is 0.823 (fine-only).
\item
  \textbf{The abstract says ``Selecting every choice on five
  class-stratified folds''} (p.~1, \texttt{sec/0\_abstract.tex:2}),
  which has the same problem.
\end{enumerate}

\subsection{1.3 How to handle it carefully}

\begin{enumerate}
\def\labelenumi{\arabic{enumi}.}
\tightlist
\item
  \textbf{Say plainly what the rule chose and why} (sheep guard; uniform
  weights), next to our choice. Print both pipelines on both benchmarks
  (the rule's choice is not printed anywhere in the main paper today;
  see \#2).
\item
  \textbf{Say when the override happened}: after COCO-Cap was seen.
  Treat the COCO-Cap gain as evidence for the calibration analysis (Sec.
  5.5), not as held-out confirmation.
\item
  \textbf{Label 0.829 as forced OOF} everywhere.
\item
  \textbf{Reword contribution 4}: a protocol that \emph{records} every
  decision, including the two owner decisions, rather than one that
  ``made every choice''.
\item
  \textbf{Give a clean test of the decision}, scored after it was
  frozen. We now have one: the \textbf{\#5 backbone runs} (Sec. \#5
  below) were done after the hierarchy was frozen, on new backbones. The
  hierarchy holds on SD 2.1 (+0.054 / +0.060) and fails on SDXL (\ensuremath{-}0.022
  / \ensuremath{-}0.011). That is honest, untouched evidence, positive and negative.
  (Caveat: that setting is one noising pass, not the paper's
  generation-time setting.)
\end{enumerate}

\subsection{1.4 The evidence behind each claim, and what happens to it}

{\def\LTcaptype{none} % do not increment counter
\begin{longtable}[]{@{}
  >{\raggedright\arraybackslash}p{(\linewidth - 6\tabcolsep) * \real{0.2500}}
  >{\raggedright\arraybackslash}p{(\linewidth - 6\tabcolsep) * \real{0.2500}}
  >{\raggedright\arraybackslash}p{(\linewidth - 6\tabcolsep) * \real{0.2500}}
  >{\raggedright\arraybackslash}p{(\linewidth - 6\tabcolsep) * \real{0.2500}}@{}}
\toprule\noalign{}
\begin{minipage}[b]{\linewidth}\raggedright
Claim in the paper
\end{minipage} & \begin{minipage}[b]{\linewidth}\raggedright
Where (JC's PDF; repo file)
\end{minipage} & \begin{minipage}[b]{\linewidth}\raggedright
Evidence now
\end{minipage} & \begin{minipage}[b]{\linewidth}\raggedright
Action
\end{minipage} \\
\midrule\noalign{}
\endhead
\bottomrule\noalign{}
\endlastfoot
``+0.125 / +0.114 over class-word OVAM'' & p.~8, Sec. 6;
\texttt{sec/6\_comparison.tex:3}. Abstract and intro give 0.829 / 0.696
``against 0.704 and 0.582 for class-word OVAM'' (p.~1) & Table 1 (p.~4):
new components add 0.012 of it & A row in the comparison table, not the
headline \\
``Most of the COCO-Cap margin over our fine-only pipeline comes from the
capped hierarchy'' (+0.034) & p.~8, Sec. 6; Table 1 last row (p.~4);
Table 3 (p.~7) & F's trace (workplan Entry 9): the +0.034 compares the
two maps at \textbf{two different thresholds} (0.55 vs 0.425). At one
shared threshold: +0.012 {[}\ensuremath{-}0.002, +0.026{]}. With each map at its own
COCO-best threshold: +0.018. The per-image oracle is a tie (0.896 vs
0.895) & Keep, restated as \textbf{cross-image calibration} (which is
what Sec. 5.5 already says) \\
SA-guided layer weighting helps (+0.018 on COCO-Cap) & p.~4, Table 1 row
7; p.~5--6, Sec. 5.1 & Rerun (Entry 8): \textbf{+0.024 {[}+0.013,
+0.036{]}} on COCO-Cap at 0.55 & Keep: the cleanest positive result in
clutter \\
``DenseCRF gains six times more on capped maps in cluttered scenes'' &
p.~1, abstract; p.~6, Sec. 5.4; p.~7, Table 4 caption;
\texttt{sec/0\_abstract.tex:2}, \texttt{sec/5\_analysis.tex:50},
\texttt{tables/tab\_refine.tex:15} & Same two-threshold artefact: at
matched thresholds DenseCRF adds about +0.014 to both maps (N audit) &
\textbf{Drop} \\
One-sided cap beats no cap & p.~6, Sec. 5.3; p.~7, Table 3 & Rerun:
+0.010 VOC-Sim; \textbf{\ensuremath{-}0.005 {[}\ensuremath{-}0.011, \ensuremath{-}0.001{]}} COCO-Cap at 0.55 &
Weaken: keep the mass analysis (13\% background vs 0.5\% object
removed) \\
``Every choice is made on VOC-Sim folds, with COCO-Cap used only after
freezing'' & p.~1, intro; \texttt{sec/1\_intro.tex:7} & Override after
COCO-Cap (protocol Section 135) & \textbf{Rewrite} (below) \\
Best training-free method (implied) & p.~8, Table 5 & Paired comparison
(\#3): NERVE (adds CLIP) beats BRACE on COCO-Cap, +0.032 {[}+0.007,
+0.060{]} & Narrow to ``best method that uses only the diffusion
model'' \\
\end{longtable}
}

\subsection{1.5 Replacement text, location by location}

Numbers below are the paper's current ones (0.829 / 0.696). If the
owners switch to the prompt-corrected capture (runs/20: 0.832 / 0.709;
see 1.6), replace them in one pass.

\textbf{(a) Abstract} --- p.~1; \texttt{sec/0\_abstract.tex:2} (whole
abstract).

\emph{Now:} ``\ldots{} We show that how these signals are combined
across spatial scales matters more than which graph operator is applied
to them. \ldots{} DenseCRF gains six times more on capped maps in
cluttered scenes. Selecting every choice on five class-stratified folds
of VOC-Sim, the pipeline reaches 0.829 out-of-fold mIoU and transfers
frozen to COCO-Cap at 0.696, against 0.704 and 0.582 for class-word
OVAM. \ldots{}''

\emph{Replace with:}

\begin{quote}
Text-to-image diffusion models localize words in cross-attention (CA)
and group pixels in self-attention (SA), so the images they generate can
be segmented without training. We study how these signals should be
combined across spatial scales, and find that the choice matters most in
cluttered scenes. Our training-free pipeline, BRACE, weights CA layers
by their agreement with SA, propagates the resulting class-word seed
from coarse to fine SA, and lets each finer level redistribute evidence
only within the mass its coarse parent assigns. All settings are
selected on class-stratified folds of the single-object VOC-Sim, where
the two new components add little (+0.006 each; 0.829 forced out-of-fold
mIoU). On the cluttered COCO-Cap the frozen pipeline reaches 0.696,
0.034 above a fine-only pipeline selected the same way, and SA-guided
layer weighting is the most robust single component. An oracle analysis
shows why: the hierarchy does not rank pixels better within an image; it
makes maps comparable across images, so a threshold chosen on simple
scenes still works on cluttered ones. In a paired comparison on the same
images, BRACE is the strongest method we evaluated that uses only the
diffusion model; a method that adds a CLIP seed is stronger on COCO-Cap,
and we report it separately. Our registered selection rule preferred the
simpler fine-only pipeline; we adopted the hierarchy after observing
COCO-Cap, and report both pipelines and every negative result.
\end{quote}

\textbf{(b) Introduction, results paragraph} --- p.~1, right column
(``The resulting pipeline, BRACE \ldots{}'');
\texttt{sec/1\_intro.tex:7}.

\emph{Now:} ``\ldots{} On VOC-Sim it reaches 0.829 out-of-fold mIoU, and
on COCO-Cap it reaches 0.696 in frozen transfer, against 0.704 and 0.582
for OVAM with class words. Every choice is made on VOC-Sim folds, with
COCO-Cap used only after freezing.''

\emph{Replace with:}

\begin{quote}
The resulting pipeline, BRACE (Bounded Refinement of Attention from
Coarse Envelopes), uses no training, no optimized tokens and no tuned
graph parameters. Its two new components add little on single objects
(0.829 forced out-of-fold mIoU on VOC-Sim, +0.006 each) but matter in
clutter: on COCO-Cap the hierarchy adds 0.034 over fine-only
propagation, and we show that this gain is cross-image calibration (Sec.
5.5). Every setting is selected on VOC-Sim folds; the decision to adopt
the hierarchy over the simpler pipeline preferred by our registered rule
was made after COCO-Cap had been scored, so we report both pipelines
throughout.
\end{quote}

\textbf{(c) Contribution list} --- p.~1--2;
\texttt{sec/1\_intro.tex:11--14}.

\emph{Replace with:}

\begin{quote}
\begin{enumerate}
\def\labelenumi{\arabic{enumi}.}
\tightlist
\item
  A diffusion-only, training-free pipeline that turns attention into
  class masks: SA-guided CA layer weighting, coarse-to-fine SA
  propagation, and a one-sided cap that bounds each finer level by the
  coarser one.
\item
  Evidence that the scale structure matters in clutter rather than on
  single objects: about +0.006 per component on VOC-Sim, +0.034 for the
  hierarchy and +0.018 for the layer weighting on COCO-Cap, with
  controlled ablations of each choice.
\item
  An oracle threshold analysis that explains the hierarchy's gain as
  cross-image calibration, not better ranking within an image, and its
  main failure (objects that fill the frame).
\item
  A fold-safe protocol that records every decision --- including the two
  components adopted against the registered rule, and when --- and
  reports all negative results, together with paired reruns of competing
  methods on the same images.
\end{enumerate}
\end{quote}

\textbf{(d) Sec. 4, ``Selection status''} --- p.~4;
\texttt{sec/4\_protocol.tex:10}.

\emph{Now:} ``\ldots{} Two components were adopted after the rule had
retained a simpler reference. The self-attention-guided layer weights
won 4/5 folds and transferred, but the multi-candidate screen kept
uniform weights. The capped hierarchy won 5/5 folds against both
references, but one class (sheep, Sec. 5.5) lost more than the 0.05
guard. The supplementary material reports the registered outcomes in
full.''

\emph{Replace with:}

\begin{quote}
The CA scales, the SA module, the two-pass baseline, DenseCRF and the SA
call window were promoted by this rule. The rule kept simpler references
for the two new components. The SA-guided layer weights won 4/5 folds,
but the multi-candidate screen kept uniform weights. The capped
hierarchy won 5/5 folds against both references (+0.006 forced OOF), but
one class (sheep, Sec. 5.5) lost more than the 0.05 guard, so the rule
kept fine-only propagation. We adopted both components as an explicit
decision after the COCO-Cap results had been computed. COCO-Cap is
therefore not an untouched test set for these two decisions, although
every setting inside them was selected on VOC-Sim. We report the rule's
pipeline next to ours on both benchmarks (Table X), use the COCO-Cap
gain as evidence for the calibration analysis of Sec. 5.5 rather than as
held-out confirmation, and give a test scored after the decision was
frozen (new backbones, Sec. 5.y). Section S8 and Tab. S10 list the
registered outcomes.
\end{quote}

\textbf{(e) Sec. 5.4 and Table 4 caption, ``six times''} --- p.~6 (Sec.
5.4) and p.~7 (Table 4); \texttt{sec/5\_analysis.tex:50},
\texttt{tables/tab\_refine.tex:15}.

\emph{Now (both places):} ``\ldots{} On capped coarse-to-fine maps it
raises precision, helps COCO-Cap six times more (+0.014, 18 of 20
classes), \ldots{}''

\emph{Replace with:}

\begin{quote}
On capped coarse-to-fine maps it mainly raises precision and improves
427 COCO-Cap images against 105 that get worse (fine-only: 354 against
243).
\end{quote}

and delete ``DenseCRF gains six times more on capped maps in cluttered
scenes.'' from the abstract (done in (a)).

\textbf{(f) Sec. 6 (will become Sec. 5.1, see \#8)} --- p.~8;
\texttt{sec/6\_comparison.tex:3}.

\emph{Now:} ``\ldots{} With class words only and no training, it exceeds
OVAM by 0.125 on VOC-Sim and 0.114 on COCO-Cap. It also exceeds
DatasetDM, which trains an extra decoder, on both benchmarks. Most of
the COCO-Cap margin over our fine-only pipeline comes from the capped
hierarchy: cluttered scenes benefit most from a coarse envelope.
\ldots{}''

\emph{Replace with:}

\begin{quote}
Table 5 compares BRACE with competing methods on the same images (paired
block) and lists the numbers published for these benchmarks (unpaired
block). Among methods that use only the diffusion model, BRACE is the
strongest on both benchmarks: GoCA trails it by 0.08 / 0.11 with SD 1.5
and by 0.03 / 0.10 with SDXL, and DiffSegmenter by 0.04 / 0.02. NERVE,
which adds a CLIP seed, is on par on VOC-Sim and stronger on COCO-Cap
(+0.032 {[}+0.007, +0.060{]}); with its own released thresholding rule
it falls below BRACE (0.812 / 0.663). The published numbers come from
OVAM's reruns of each method and use other images and post-processing,
so they are indicative only.
\end{quote}

(OVAM, DAAM and SCLIP rows {[}pending{]}, see \#3.)

\textbf{(g) Sec. 7, ``Failure modes and scope'', last sentence} ---
p.~8; \texttt{sec/7\_discussion.tex:7}.

\emph{Now:} ``Two components were adopted after the registered rule kept
a simpler reference (Sec. 4); their gains hold in 4/5 and 5/5 held-out
folds and transfer, but they are not pre-registered promotions.''

\emph{Replace with:}

\begin{quote}
Two components were adopted after the registered rule kept a simpler
reference, and after COCO-Cap had been scored (Sec. 4). Their gains hold
in 4/5 and 5/5 held-out folds; on new backbones, scored after the
decision was frozen, the hierarchy helps on SD 2.1 but not on SDXL (Sec.
5.y). A further benchmark scored once with both pipelines frozen would
test the decision cleanly.
\end{quote}

\textbf{(h) Conclusion} --- p.~8; \texttt{sec/7\_discussion.tex:10}. Add
one sentence:

\begin{quote}
The division of labour matters most in cluttered scenes, where it makes
maps comparable across images, and it needs at least three attention
scales: on SDXL's two, the hierarchy does not help.
\end{quote}

\subsection{1.6 Open decisions for the owners}

\begin{itemize}
\tightlist
\item
  \textbf{Which capture is canonical}: the original (0.829 / 0.696) or
  the prompt-corrected one (runs/19--20: some benchmark prompts did not
  match the stored images; BRACE 0.832 / 0.709). All new paired
  comparisons (\#3, \#5) use the corrected prompts.
\item
  \textbf{Headline}: ``best diffusion-only method + separate CLIP
  block'' (recommended) or an analysis paper without a ranking claim.
\end{itemize}

\textbf{Slides:} \texttt{c1\_N.tex} (5 frames: claim moves; claims kept;
claims weakened or dropped; override disclosure; options).

\begin{center}\rule{0.5\linewidth}{0.5pt}\end{center}

\section{\#2 Mapping contributions to evidence}

\begin{quote}
\emph{JC:} ``for each of the four contributions \ldots{} which section,
table and figure provide the evidence, and which number is the key
result.''
\end{quote}

\subsection{2.1 The mapping (current paper, JC's PDF numbering; ``S'' =
supplementary)}

{\def\LTcaptype{none} % do not increment counter
\begin{longtable}[]{@{}
  >{\raggedright\arraybackslash}p{(\linewidth - 12\tabcolsep) * \real{0.1429}}
  >{\raggedright\arraybackslash}p{(\linewidth - 12\tabcolsep) * \real{0.1429}}
  >{\raggedright\arraybackslash}p{(\linewidth - 12\tabcolsep) * \real{0.1429}}
  >{\raggedright\arraybackslash}p{(\linewidth - 12\tabcolsep) * \real{0.1429}}
  >{\raggedright\arraybackslash}p{(\linewidth - 12\tabcolsep) * \real{0.1429}}
  >{\raggedright\arraybackslash}p{(\linewidth - 12\tabcolsep) * \real{0.1429}}
  >{\raggedright\arraybackslash}p{(\linewidth - 12\tabcolsep) * \real{0.1429}}@{}}
\toprule\noalign{}
\begin{minipage}[b]{\linewidth}\raggedright
\#
\end{minipage} & \begin{minipage}[b]{\linewidth}\raggedright
Contribution (Sec. 1, p.~1)
\end{minipage} & \begin{minipage}[b]{\linewidth}\raggedright
Section
\end{minipage} & \begin{minipage}[b]{\linewidth}\raggedright
Tables
\end{minipage} & \begin{minipage}[b]{\linewidth}\raggedright
Figures
\end{minipage} & \begin{minipage}[b]{\linewidth}\raggedright
Key number
\end{minipage} & \begin{minipage}[b]{\linewidth}\raggedright
Main or supp
\end{minipage} \\
\midrule\noalign{}
\endhead
\bottomrule\noalign{}
\endlastfoot
1 & \textbf{Pipeline}: SA-guided seed, coarse-to-fine SA with cap,
DenseCRF & 3.2--3.5, 5, 6 & Tab. 1 (last row), Tab. 5; S6, S11 & Fig. 1,
2; S8, S9 & \textbf{0.829} forced OOF / \textbf{0.696} frozen; vs OVAM
0.704 / 0.582 (published, not paired) & Main \\
2a & Ablation: CA scales and layer weighting & 3.2, 3.3, 5.1 & Tab. 2,
Tab. 1 rows 1--2, 7; S1, S3 & none in main; S1, S4 & drop 32-px CA:
0.665 \ensuremath{\rightarrow} 0.691 (4/5 folds); weighting 0.817 \ensuremath{\rightarrow} 0.823 (4/5), COCO 0.644 \ensuremath{\rightarrow}
0.662 & Numbers main; \textbf{pictures only in supp} \\
2b & Ablation: coarse vs fine SA & 3.4, 5.2 & Tab. 3 (first rows); S4 &
Fig. 3 left/middle & 16-px SA alone 0.72; hierarchy vs fine-only COCO
0.693 vs 0.662 & Main; \textbf{no picture of a fine-only map next to the
coarse one} \\
2c & Ablation: cap vs exact conservation vs none & 3.4, 5.3 & Tab. 3; S5
& Fig. 3 right, Fig. 4; S8 & cap 0.830 (5/5 folds); conservation \ensuremath{-}5
points raw; cap removes 13\% background vs 0.5\% object & Main \\
2d & Ablation: DenseCRF on each map & 3.5, 5.4 & Tab. 4, Tab. 1 row 5 &
Fig. 5; S5 & images up/down 427/105 (capped) vs 354/243 (fine-only) on
COCO-Cap (``six times'' to be dropped, \#1) & Main \\
3 & \textbf{Oracle threshold analysis} & 5.5, 7 & \textbf{none in main};
S7 & Fig. 6; S6, S7 & per-image oracle tie 0.896 vs 0.895; best global
threshold 0.425 \ensuremath{\rightarrow} 0.55; sheep \ensuremath{-}0.09 & Claim main; \textbf{support in
supp} \\
4 & \textbf{Fold-safe protocol} & 4, 7 & \textbf{none in main}; S2, S8,
S9, S10 & \textbf{none} & rule (+0.010 / 0.005, 12/20, 0.05 guard, 4/5
folds); 45 registered comparisons; \textbf{the rule's own choice
(fine-only, 0.823) is not printed in the main paper} & \textbf{Mostly
supp} \\
\end{longtable}
}

Full version with the committed source of every number:
\texttt{comparisons/notes/entry\_2\_mapping.md}.

\subsection{2.2 What the table shows}

\begin{itemize}
\tightlist
\item
  \textbf{Well supported in the main paper}: contributions 1 and 2 (all
  four ablations).
\item
  \textbf{Rely on the supplementary}: 3 (the oracle has a figure but its
  supporting table is S7) and \textbf{4} (no table or figure at all in
  the main paper).
\item
  \textbf{Weak spots}: the qualitative claim of contribution 2 (``each
  is backed by quantitative and qualitative evidence'') holds for the
  cap and DenseCRF, not for CA scales (pictures only in supp) or coarse
  vs fine SA (no fine-only picture).
\item
  \textbf{Results that support no contribution} (shorten or move to the
  supplement): Sec. 5.1's untracked sentence on CA heads and calls; the
  COCO-Cap person column in Fig. 2; the GoCA-orientation footnote and
  Tab. S1; Tab. S9 (our adaptations of related methods, now superseded
  by real reruns, \#3); Sec. S9 / Tab. S11 (merge into the new Table 5);
  Sec. S1 ``Cost'' (replaced by measured cost, \#8).
\end{itemize}

\subsection{2.3 Fixes that follow}

\begin{itemize}
\tightlist
\item
  Contribution 4: add a small main-paper table, \textbf{the rule's
  choice next to ours} (fine-only + uniform weights vs BRACE, both
  benchmarks). This also handles \#1.
\item
  Contribution 3: move one row of Tab. S7 (the label-free per-image
  thresholds) into the text of Sec. 5.5.
\end{itemize}

\textbf{Slides:} \texttt{c2\_N.tex}.

\begin{center}\rule{0.5\linewidth}{0.5pt}\end{center}

\section{\#3 Comparison with the state of the art}

\begin{quote}
\emph{JC:} ``a paired block on the same images: our own OVAM
reimplementation as the baseline (also reporting OVAM's optimized-token
results, flagged as a different setting), GoCA, \ldots{} and if possible
NERVE \ldots{} PEARL \ldots{} cite it for positioning \ldots{} The
target Table 5 would have two blocks: published numbers (marked as not
paired) and the paired comparison.''
\end{quote}

\subsection{3.1 What we did and how (configuration)}

{\def\LTcaptype{none} % do not increment counter
\begin{longtable}[]{@{}
  >{\raggedright\arraybackslash}p{(\linewidth - 2\tabcolsep) * \real{0.5000}}
  >{\raggedright\arraybackslash}p{(\linewidth - 2\tabcolsep) * \real{0.5000}}@{}}
\toprule\noalign{}
\begin{minipage}[b]{\linewidth}\raggedright
Setting
\end{minipage} & \begin{minipage}[b]{\linewidth}\raggedright
Value
\end{minipage} \\
\midrule\noalign{}
\endhead
\bottomrule\noalign{}
\endlastfoot
Images & the same 490 VOC-Sim + 707 COCO-Cap images and annotations for
every method (ground truth checked identical on 1,197/1,197) \\
Query & the class word only, for every method \\
Code & each method from its own public repository and settings (GoCA
github.com/Darkbblue/goca; NERVE; DiffSegmenter); OVAM and DAAM from
their released code {[}pending{]} \\
BRACE & regenerated with the prompt-corrected capture (runs/20), F's
code imported unchanged; reproduces runs/20 to 0.001 \\
Threshold & chosen on VOC-Sim (5-fold out-of-fold), reused unchanged on
COCO-Cap; also fixed 0.5 and each method's own rule (in
\texttt{result.md}) \\
Post-processing & the same DenseCRF (OVAM's parameters) for every
method; raw numbers also reported \\
Statistics & paired \ensuremath{\Delta} (method \ensuremath{-} BRACE), 2,000-sample image bootstrap,
95\% interval \\
Backbone & SD 1.5 (GoCA also SD 2.1-base and SDXL) \\
Files & \texttt{comparisons/official/result.md} (pushed, workplan
Entries 11 and 14) \\
\end{longtable}
}

\textbf{Which benchmarks the methods report} (checked in each paper):
GoCA, NERVE, DiffSegmenter, SPARK and PEARL report only real-image
benchmarks; none reports VOC-Sim or COCO-Cap. So they enter only through
our paired reruns. SPARK has no public code and takes no class query:
no-go. PEARL (CLIP/DINO, real images, multi-class): related work only.

\textbf{The published block is OVAM's reruns.} DAAM, Attn2Mask, Grounded
Diffusion and DatasetDM did not publish VOC-Sim / COCO-Cap numbers
themselves; Table 5's rows are OVAM's runs of them (the Table 5 caption
says so; supp. line 434 says otherwise and must be fixed). Rescoring
Attn2Mask's released masks gives 0.603 on VOC-Sim against OVAM's printed
.687 (unexplained, to check).

\subsection{3.2 The new Table 5}

VOC-Sim / COCO-Cap mIoU. Paired rows: DenseCRF, VOC-selected threshold.
\ensuremath{\Delta} = method \ensuremath{-} BRACE.

{\def\LTcaptype{none} % do not increment counter
\begin{longtable}[]{@{}
  >{\raggedright\arraybackslash}p{(\linewidth - 12\tabcolsep) * \real{0.1429}}
  >{\raggedright\arraybackslash}p{(\linewidth - 12\tabcolsep) * \real{0.1429}}
  >{\raggedright\arraybackslash}p{(\linewidth - 12\tabcolsep) * \real{0.1429}}
  >{\raggedright\arraybackslash}p{(\linewidth - 12\tabcolsep) * \real{0.1429}}
  >{\raggedright\arraybackslash}p{(\linewidth - 12\tabcolsep) * \real{0.1429}}
  >{\raggedright\arraybackslash}p{(\linewidth - 12\tabcolsep) * \real{0.1429}}
  >{\raggedright\arraybackslash}p{(\linewidth - 12\tabcolsep) * \real{0.1429}}@{}}
\toprule\noalign{}
\begin{minipage}[b]{\linewidth}\raggedright
Block
\end{minipage} & \begin{minipage}[b]{\linewidth}\raggedright
Method
\end{minipage} & \begin{minipage}[b]{\linewidth}\raggedright
Extra model
\end{minipage} & \begin{minipage}[b]{\linewidth}\raggedright
VOC-Sim
\end{minipage} & \begin{minipage}[b]{\linewidth}\raggedright
COCO-Cap
\end{minipage} & \begin{minipage}[b]{\linewidth}\raggedright
\ensuremath{\Delta} VOC-Sim
\end{minipage} & \begin{minipage}[b]{\linewidth}\raggedright
\ensuremath{\Delta} COCO-Cap
\end{minipage} \\
\midrule\noalign{}
\endhead
\bottomrule\noalign{}
\endlastfoot
\textbf{Paired, diffusion only, class word} & \textbf{BRACE (ours)} &
--- & \textbf{0.831} & \textbf{0.708} & --- & --- \\
& fine-only (ours, the rule's choice) & --- & {[}pending{]} &
{[}pending{]} & {[}pending{]} & {[}pending{]} \\
& OVAM, class word (our run) & --- & {[}pending{]} & {[}pending{]} &
{[}pending{]} & {[}pending{]} \\
& DAAM (our run) & --- & {[}pending{]} & {[}pending{]} & {[}pending{]} &
{[}pending{]} \\
& DiffSegmenter & --- & 0.795 & 0.688 & \ensuremath{-}0.036 {[}\ensuremath{-}0.049, \ensuremath{-}0.023{]} &
\ensuremath{-}0.020 {[}\ensuremath{-}0.041, +0.001{]} \\
& GoCA, SD 1.5 & --- & 0.753 & 0.599 & \ensuremath{-}0.078 {[}\ensuremath{-}0.095, \ensuremath{-}0.060{]} &
\ensuremath{-}0.110 {[}\ensuremath{-}0.134, \ensuremath{-}0.084{]} \\
& GoCA, SD 2.1-base & --- & 0.740 & 0.597 & \ensuremath{-}0.091 {[}\ensuremath{-}0.110, \ensuremath{-}0.072{]}
& \ensuremath{-}0.111 {[}\ensuremath{-}0.134, \ensuremath{-}0.086{]} \\
& GoCA, SDXL & --- & 0.797 & 0.606 & \ensuremath{-}0.034 {[}\ensuremath{-}0.048, \ensuremath{-}0.021{]} &
\ensuremath{-}0.102 {[}\ensuremath{-}0.124, \ensuremath{-}0.078{]} \\
\textbf{Paired, with an extra vision-language model} & NERVE, class word
& CLIP ViT-L/14 & 0.844 & 0.740 & +0.012 {[}\ensuremath{-}0.002, +0.028{]} &
\textbf{+0.032 {[}+0.007, +0.060{]}} \\
& NERVE, 20-class vocabulary & CLIP ViT-L/14 & 0.839 & 0.747 & +0.008
{[}\ensuremath{-}0.007, +0.023{]} & \textbf{+0.039 {[}+0.013, +0.068{]}} \\
& SCLIP + DenseCRF & CLIP ViT-B/16 & {[}pending{]} & {[}pending{]} &
{[}pending{]} & {[}pending{]} \\
& SCLIP seed + BRACE hierarchy (F, runs/20) & CLIP ViT-B/16 & 0.852 &
0.734 & & \\
\textbf{Paired, few-shot (optimized token; not training-free)} & OVAM,
optimized token (our run) & --- & {[}pending{]} & {[}pending{]} &
{[}pending{]} & {[}pending{]} \\
& BRACE, optimized token (\#4) & --- & {[}pending{]} (Entry 8, old
prompts: 0.861 / 0.681 at 0.55) & & & \\
\textbf{Published, not paired} (OVAM's Table 1; other images and
post-processing) & DAAM & --- & .662 & .484 & & \\
& Attn2Mask & --- & .687 & .550 & & \\
& OVAM, class word & --- & .704 & .582 & & \\
& OVAM, optimized token & --- & .825 & --- & & \\
& Grounded Diffusion (trains a module) & --- & .621 & .502 & & \\
& DatasetDM (trains a decoder) & --- & .803 & .593 & & \\
\end{longtable}
}

\textbf{Reading.} BRACE is the best method that uses only the diffusion
model. NERVE, which adds CLIP, is tied on VOC-Sim and better on
COCO-Cap. With its own released thresholding rule NERVE drops to 0.812 /
0.663, below BRACE. Without DenseCRF every sign is the same.

\textbf{Caveats to print in the caption.} BRACE reads attention while
generating, so it regenerates each image; the other methods segment the
stored image (the two differ by under one grey level on most images;
every sign holds on the closest subset). DiffSegmenter's own evaluation
tunes its background score on the test set, so it has no fair own-rule
row. All rows are exploratory; BRACE's VOC-Sim number is in-sample for
its design.

\textbf{LaTeX draft} (replaces \texttt{tables/tab\_published.tex}; fill
{[}pending{]} rows when done):

\begin{Shaded}
\begin{Highlighting}[]
\KeywordTok{\textbackslash{}begin}\NormalTok{\{}\ExtensionTok{table}\NormalTok{\}[t]}
\FunctionTok{\textbackslash{}centering}
\FunctionTok{\textbackslash{}footnotesize\textbackslash{}setlength}\NormalTok{\{}\FunctionTok{\textbackslash{}tabcolsep}\NormalTok{\}\{3pt\}}
\KeywordTok{\textbackslash{}begin}\NormalTok{\{}\ExtensionTok{tabular}\NormalTok{\}\{}\CharTok{@\{\}}\NormalTok{lcccc}\CharTok{@\{\}}\NormalTok{\}}
\FunctionTok{\textbackslash{}toprule}
\NormalTok{Method }\OperatorTok{\&}\NormalTok{ VOC{-}Sim }\OperatorTok{\&}\NormalTok{ COCO{-}Cap }\OperatorTok{\&} \SpecialStringTok{$}\SpecialCharTok{\textbackslash{}Delta}\SpecialStringTok{$}\NormalTok{ VOC }\OperatorTok{\&} \SpecialStringTok{$}\SpecialCharTok{\textbackslash{}Delta}\SpecialStringTok{$}\NormalTok{ COCO }\FunctionTok{\textbackslash{}\textbackslash{}}
\FunctionTok{\textbackslash{}midrule}
\FunctionTok{\textbackslash{}multicolumn}\NormalTok{\{5\}\{}\CharTok{@\{\}}\NormalTok{l\}\{}\FunctionTok{\textbackslash{}emph}\NormalTok{\{Paired, same images, diffusion only, class word\}\}}\FunctionTok{\textbackslash{}\textbackslash{}}
\FunctionTok{\textbackslash{}textbf}\NormalTok{\{BRACE (ours)\} }\OperatorTok{\&} \FunctionTok{\textbackslash{}textbf}\NormalTok{\{.831\} }\OperatorTok{\&} \FunctionTok{\textbackslash{}textbf}\NormalTok{\{.708\} }\OperatorTok{\&}\NormalTok{ {-}{-} }\OperatorTok{\&}\NormalTok{ {-}{-} }\FunctionTok{\textbackslash{}\textbackslash{}}
\NormalTok{DiffSegmenter\textasciitilde{}}\KeywordTok{\textbackslash{}cite}\NormalTok{\{}\ExtensionTok{...}\NormalTok{\} }\OperatorTok{\&}\NormalTok{ .795 }\OperatorTok{\&}\NormalTok{ .688 }\OperatorTok{\&} \SpecialStringTok{${-}$}\NormalTok{.036 }\OperatorTok{\&} \SpecialStringTok{${-}$}\NormalTok{.020 }\FunctionTok{\textbackslash{}\textbackslash{}}
\NormalTok{GoCA, SD 1.5\textasciitilde{}}\KeywordTok{\textbackslash{}cite}\NormalTok{\{}\ExtensionTok{meng2026making}\NormalTok{\} }\OperatorTok{\&}\NormalTok{ .753 }\OperatorTok{\&}\NormalTok{ .599 }\OperatorTok{\&} \SpecialStringTok{${-}$}\NormalTok{.078 }\OperatorTok{\&} \SpecialStringTok{${-}$}\NormalTok{.110 }\FunctionTok{\textbackslash{}\textbackslash{}}
\NormalTok{GoCA, SDXL\textasciitilde{}}\KeywordTok{\textbackslash{}cite}\NormalTok{\{}\ExtensionTok{meng2026making}\NormalTok{\} }\OperatorTok{\&}\NormalTok{ .797 }\OperatorTok{\&}\NormalTok{ .606 }\OperatorTok{\&} \SpecialStringTok{${-}$}\NormalTok{.034 }\OperatorTok{\&} \SpecialStringTok{${-}$}\NormalTok{.102 }\FunctionTok{\textbackslash{}\textbackslash{}}
\FunctionTok{\textbackslash{}midrule}
\FunctionTok{\textbackslash{}multicolumn}\NormalTok{\{5\}\{}\CharTok{@\{\}}\NormalTok{l\}\{}\FunctionTok{\textbackslash{}emph}\NormalTok{\{Paired, with a CLIP model\}\}}\FunctionTok{\textbackslash{}\textbackslash{}}
\NormalTok{NERVE\textasciitilde{}}\KeywordTok{\textbackslash{}cite}\NormalTok{\{}\ExtensionTok{mahatha2025nerve}\NormalTok{\} }\OperatorTok{\&}\NormalTok{ .844 }\OperatorTok{\&}\NormalTok{ .740 }\OperatorTok{\&}\NormalTok{ +.012 }\OperatorTok{\&} \FunctionTok{\textbackslash{}textbf}\NormalTok{\{+.032\} }\FunctionTok{\textbackslash{}\textbackslash{}}
\FunctionTok{\textbackslash{}midrule}
\FunctionTok{\textbackslash{}multicolumn}\NormalTok{\{5\}\{}\CharTok{@\{\}}\NormalTok{l\}\{}\FunctionTok{\textbackslash{}emph}\NormalTok{\{Published (OVAM\textquotesingle{}s runs; not paired)\}\}}\FunctionTok{\textbackslash{}\textbackslash{}}
\NormalTok{DAAM\textasciitilde{}}\KeywordTok{\textbackslash{}cite}\NormalTok{\{}\ExtensionTok{tang2023daam}\NormalTok{\} }\OperatorTok{\&}\NormalTok{ .662 }\OperatorTok{\&}\NormalTok{ .484 }\OperatorTok{\&} \OperatorTok{\&} \FunctionTok{\textbackslash{}\textbackslash{}}
\NormalTok{Attn2Mask\textasciitilde{}}\KeywordTok{\textbackslash{}cite}\NormalTok{\{}\ExtensionTok{yoshihashi2023attention}\NormalTok{\} }\OperatorTok{\&}\NormalTok{ .687 }\OperatorTok{\&}\NormalTok{ .550 }\OperatorTok{\&} \OperatorTok{\&} \FunctionTok{\textbackslash{}\textbackslash{}}
\NormalTok{OVAM\textasciitilde{}}\KeywordTok{\textbackslash{}cite}\NormalTok{\{}\ExtensionTok{marcosmanchon2024ovam}\NormalTok{\} }\OperatorTok{\&}\NormalTok{ .704 }\OperatorTok{\&}\NormalTok{ .582 }\OperatorTok{\&} \OperatorTok{\&} \FunctionTok{\textbackslash{}\textbackslash{}}
\NormalTok{Grounded Diffusion\textasciitilde{}}\KeywordTok{\textbackslash{}cite}\NormalTok{\{}\ExtensionTok{li2023open}\NormalTok{\} }\OperatorTok{\&}\NormalTok{ .621 }\OperatorTok{\&}\NormalTok{ .502 }\OperatorTok{\&} \OperatorTok{\&} \FunctionTok{\textbackslash{}\textbackslash{}}
\NormalTok{DatasetDM\textasciitilde{}}\KeywordTok{\textbackslash{}cite}\NormalTok{\{}\ExtensionTok{wu2023datasetdm}\NormalTok{\} }\OperatorTok{\&}\NormalTok{ .803 }\OperatorTok{\&}\NormalTok{ .593 }\OperatorTok{\&} \OperatorTok{\&} \FunctionTok{\textbackslash{}\textbackslash{}}
\FunctionTok{\textbackslash{}bottomrule}
\KeywordTok{\textbackslash{}end}\NormalTok{\{}\ExtensionTok{tabular}\NormalTok{\}}
\FunctionTok{\textbackslash{}caption}\NormalTok{\{}\FunctionTok{\textbackslash{}textbf}\NormalTok{\{Comparison with competing methods.\} Paired rows:}
\NormalTok{  the same 490 VOC{-}Sim and 707 COCO{-}Cap images, class word, DenseCRF,}
\NormalTok{  threshold selected on VOC{-}Sim and reused on COCO{-}Cap;}
  \SpecialStringTok{$}\SpecialCharTok{\textbackslash{}Delta}\SpecialStringTok{$}\NormalTok{ = method }\SpecialStringTok{${-}$}\NormalTok{ BRACE, bold when the 95}\FunctionTok{\textbackslash{}\%}\NormalTok{ bootstrap interval}
\NormalTok{  excludes zero. Published rows are OVAM\textquotesingle{}s reruns on other images}
\NormalTok{  and post{-}processing.\}}
\KeywordTok{\textbackslash{}label}\NormalTok{\{}\ExtensionTok{tab:published}\NormalTok{\}}
\KeywordTok{\textbackslash{}end}\NormalTok{\{}\ExtensionTok{table}\NormalTok{\}}
\end{Highlighting}
\end{Shaded}

\textbf{Slides:} \texttt{c3\_N.tex} (5 frames: methods and go/no-go;
sanity check; the paired table; paired \ensuremath{\Delta}; caveats and actions).

\begin{center}\rule{0.5\linewidth}{0.5pt}\end{center}

\section{\#4 BRACE with optimized tokens}

\begin{quote}
\emph{JC:} ``Please run BRACE with OVAM's optimized tokens \ldots{}
report it as a separate \guillemotleft{}optimized token\guillemotright{} block \ldots{} Check on which
images the released tokens were optimized \ldots{} state clearly that
this block is no longer training-free.''
\end{quote}

\subsection{4.1 Which images the tokens saw (N, Entry 5) --- done}

\begin{itemize}
\tightlist
\item
  OVAM released 57 per-image tokens (2--3 per class) optimized on
  \textbf{63 separate generated training images}
  (\texttt{voc\_sim\_train/}); the class token is their mean.
\item
  \textbf{No overlap} with VOC-Sim (no shared class + seed among 500
  rows; no identical pixels) or COCO-Cap (no shared seed; different
  prompts).
\item
  So the block is \textbf{few-shot (57 annotated images), not
  training-free}, with no test leakage.
\end{itemize}

\subsection{4.2 BRACE with the token (F, Entry 8 --- original prompts)
--- done}

Configuration: one live generation per image, both the class word and
OVAM's VOC-Sim token read from the same generation; threshold fixed at
0.55; DenseCRF; paired; pre-registered
(\texttt{runs/15\_token\_ladder}).

{\def\LTcaptype{none} % do not increment counter
\begin{longtable}[]{@{}lll@{}}
\toprule\noalign{}
& VOC-Sim & COCO-Cap \\
\midrule\noalign{}
\endhead
\bottomrule\noalign{}
\endlastfoot
BRACE, class word & 0.830 & 0.696 \\
BRACE, optimized token & \textbf{0.861} & 0.681 \\
token \ensuremath{-} word & \textbf{+0.031 {[}+0.022, +0.039{]}} & \ensuremath{-}0.015 {[}\ensuremath{-}0.030,
+0.000{]} \\
seed alone: token vs word & 0.588 \ensuremath{\rightarrow} 0.766 & 0.444 \ensuremath{\rightarrow} 0.625 \\
\end{longtable}
}

Do our components still help with the token? (\ensuremath{\Delta} with the token; with the
class word in brackets)

{\def\LTcaptype{none} % do not increment counter
\begin{longtable}[]{@{}
  >{\raggedright\arraybackslash}p{(\linewidth - 4\tabcolsep) * \real{0.3333}}
  >{\raggedright\arraybackslash}p{(\linewidth - 4\tabcolsep) * \real{0.3333}}
  >{\raggedright\arraybackslash}p{(\linewidth - 4\tabcolsep) * \real{0.3333}}@{}}
\toprule\noalign{}
\begin{minipage}[b]{\linewidth}\raggedright
Component
\end{minipage} & \begin{minipage}[b]{\linewidth}\raggedright
VOC-Sim
\end{minipage} & \begin{minipage}[b]{\linewidth}\raggedright
COCO-Cap
\end{minipage} \\
\midrule\noalign{}
\endhead
\bottomrule\noalign{}
\endlastfoot
hierarchy vs fine-only & +0.010 (+0.034) & \textbf{\ensuremath{-}0.017 {[}\ensuremath{-}0.028,
\ensuremath{-}0.007{]}} (+0.012) \\
SA-guided vs uniform weights & +0.003 (+0.004) & +0.004
(\textbf{+0.024}) \\
cap vs no cap & +0.012 (+0.010) & \ensuremath{-}0.002 (\ensuremath{-}0.005) \\
\end{longtable}
}

\textbf{Reading.} The token makes the seed much better (+0.18). With
such a strong seed the hierarchy is no longer needed; on COCO-Cap it
even hurts. So JC's ``stronger message'' (BRACE helps regardless of the
query) \textbf{does not hold}; the ``why'' is that BRACE fixes a weak
class-word seed, and a strong seed leaves nothing to fix. The COCO-Cap
drop is concentrated in a few classes (sofa 0.70 \ensuremath{\rightarrow} 0.34, car 0.76 \ensuremath{\rightarrow}
0.66; person rises 0.28 \ensuremath{\rightarrow} 0.62).

\subsection{4.3 Rerun on the corrected prompts (runs/25, full)
{[}pending{]}}

The same ladder on the prompt-corrected capture, exactly as F
pre-registered it (\texttt{runs/25\_token\_ladder\_nr/PREREG.md});
outputs in \texttt{comparisons/token\_ladder\_nr/}. Scoring is running.

\subsection{4.4 Paper}

\begin{itemize}
\tightlist
\item
  New ``optimized token (few-shot)'' rows in Table 5 (see \#3) and a
  short paragraph in Sec. 5.1 (new numbering).
\item
  Limitation sentence for Sec. 7: \emph{``BRACE is designed to repair a
  weak class-word seed; with OVAM's optimized token, which raises the
  seed by 0.18, the hierarchy no longer helps on COCO-Cap (\ensuremath{-}0.017). The
  tokens are optimized on 57 annotated images disjoint from both
  benchmarks, so this setting is few-shot, not training-free.''}
\end{itemize}

\textbf{Slides:} \texttt{c4\_N.tex} (provenance), F's \texttt{c4\_F.tex}
(results).

\begin{center}\rule{0.5\linewidth}{0.5pt}\end{center}

\section{\#5 Diffusion backbones}

\begin{quote}
\emph{JC:} ``evaluate fine-only, BRACE and GoCA with SD 1.5, SD 2.1 and
SDXL on the existing VOC-Sim/COCO-Cap images, adding noise and reading
the attention in one denoising pass \ldots{} SDXL has only two attention
scales, so it tests the coarse-to-fine principle itself. SD3/FLUX
\ldots{} out of scope \ldots{} If the gain does not hold, better to know
it now.''
\end{quote}

\subsection{\texorpdfstring{5.1 Configuration (pre-registered before any
capture:
\texttt{comparisons/backbones/PREREG.md})}{5.1 Configuration (pre-registered before any capture: comparisons/backbones/PREREG.md)}}

{\def\LTcaptype{none} % do not increment counter
\begin{longtable}[]{@{}
  >{\raggedright\arraybackslash}p{(\linewidth - 2\tabcolsep) * \real{0.5000}}
  >{\raggedright\arraybackslash}p{(\linewidth - 2\tabcolsep) * \real{0.5000}}@{}}
\toprule\noalign{}
\begin{minipage}[b]{\linewidth}\raggedright
Setting
\end{minipage} & \begin{minipage}[b]{\linewidth}\raggedright
Value
\end{minipage} \\
\midrule\noalign{}
\endhead
\bottomrule\noalign{}
\endlastfoot
Images & the stored benchmark images and annotations (no regeneration,
so the annotations stay valid) \\
Noise & image \ensuremath{\rightarrow} latent (VAE), Gaussian noise at timestep 100 of 1,000 (\ensuremath{\approx}
0.95 \ensuremath{\times} image + 0.33 \ensuremath{\times} noise), fixed seed 0; the same level GoCA uses \\
Pass & \textbf{one} network call per image, the image's caption as
prompt, class word as the query; no classifier-free guidance \\
Layers & SD 1.5 and SD 2.1-base: the frozen BRACE layers (identical
layout); hierarchy 16 \ensuremath{\rightarrow} 32 \ensuremath{\rightarrow} 64. SDXL: all 70 CA layers, GoCA's
reference layer, hierarchy 32 \ensuremath{\rightarrow} 64 (\textbf{one step}, its only two
scales) \\
Weights & SD 1.5; SD 2.1-base from
\texttt{sd2-community/stable-diffusion-2-1-base} (official weights
removed; hash-identical to a clone of Stability's repo); SDXL base
1.0 \\
Scoring & raw maps: threshold chosen on VOC-Sim (primary) and fixed
0.55; DenseCRF only at 0.55 (the full grid would need 240 CPU-hours) \\
GoCA & same images and masks, GoCA's own code
(\texttt{config\_model\_15/21/xl.py}, t = 100) \\
Files & \texttt{comparisons/backbones/result.md},
\texttt{comparisons/goca/PREREG\_backbones.md} (workplan Entries 14,
15) \\
\end{longtable}
}

\subsection{5.2 Results (one-pass setting; not comparable with the
paper's 0.829 / 0.696)}

Fine-only vs BRACE, raw, threshold chosen on VOC-Sim (VOC-Sim /
COCO-Cap):

{\def\LTcaptype{none} % do not increment counter
\begin{longtable}[]{@{}
  >{\raggedright\arraybackslash}p{(\linewidth - 6\tabcolsep) * \real{0.2500}}
  >{\raggedright\arraybackslash}p{(\linewidth - 6\tabcolsep) * \real{0.2500}}
  >{\raggedright\arraybackslash}p{(\linewidth - 6\tabcolsep) * \real{0.2500}}
  >{\raggedright\arraybackslash}p{(\linewidth - 6\tabcolsep) * \real{0.2500}}@{}}
\toprule\noalign{}
\begin{minipage}[b]{\linewidth}\raggedright
Backbone
\end{minipage} & \begin{minipage}[b]{\linewidth}\raggedright
Fine-only
\end{minipage} & \begin{minipage}[b]{\linewidth}\raggedright
BRACE
\end{minipage} & \begin{minipage}[b]{\linewidth}\raggedright
\ensuremath{\Delta} BRACE \ensuremath{-} fine-only
\end{minipage} \\
\midrule\noalign{}
\endhead
\bottomrule\noalign{}
\endlastfoot
SD 1.5 (reference) & 0.784 / 0.612 & 0.793 / 0.650 & +0.009 {[}+0.003,
+0.016{]} / \textbf{+0.038 {[}+0.029, +0.047{]}} \\
SD 2.1-base & 0.749 / 0.599 & 0.803 / 0.658 & \textbf{+0.054 {[}+0.044,
+0.064{]} / +0.060 {[}+0.049, +0.071{]}} \\
SDXL (two scales) & 0.761 / 0.506 & 0.739 / 0.496 & \textbf{\ensuremath{-}0.022
{[}\ensuremath{-}0.030, \ensuremath{-}0.015{]} / \ensuremath{-}0.011 {[}\ensuremath{-}0.017, \ensuremath{-}0.005{]}} \\
\end{longtable}
}

BRACE vs GoCA on the same backbone (raw, VOC-sel.):

{\def\LTcaptype{none} % do not increment counter
\begin{longtable}[]{@{}lll@{}}
\toprule\noalign{}
Backbone & GoCA & \ensuremath{\Delta} BRACE \ensuremath{-} GoCA \\
\midrule\noalign{}
\endhead
\bottomrule\noalign{}
\endlastfoot
SD 1.5 & 0.738 / 0.610 & +0.055 / +0.040 \\
SD 2.1-base & 0.719 / 0.586 & +0.083 / +0.072 \\
SDXL & 0.787 / 0.602 & \textbf{\ensuremath{-}0.048 / \ensuremath{-}0.107} \\
\end{longtable}
}

All intervals exclude zero. At one fixed threshold of 0.55, BRACE beats
fine-only on all three backbones, but that scores fine-only away from
its best threshold (the same artefact as the +0.034, \#1), so we do not
claim it.

\textbf{Reading.} The hierarchy holds on three-scale UNets (SD 1.5, SD
2.1) and fails on SDXL, where only one coarse-to-fine step is possible;
there GoCA scales with the stronger model, as it claims. This is the
answer to ``why only SD 1.5'' and an honest limitation.

\subsection{5.3 Paper}

\begin{itemize}
\tightlist
\item
  New Sec. 5.y ``Other backbones'' with the two tables above (or a
  merged one) and the setting.
\item
  Scope sentence for Sec. 7: \emph{``We study multi-resolution UNet
  diffusion models. SD3 and FLUX use one resolution and joint
  text--image attention, so they offer no coarse-to-fine hierarchy and
  are out of scope. On SDXL, with two attention scales, the hierarchy
  does not help.''}
\item
  Open: Faiz to confirm the SD 2.1 copy.
\end{itemize}

\textbf{Slides:} \texttt{c5\_N.tex} (3 frames: GoCA by backbone;
hierarchy by backbone; GoCA overtakes on SDXL).

\begin{center}\rule{0.5\linewidth}{0.5pt}\end{center}

\section{\#6 Consistency checks --- done (fix list ready)}

\begin{quote}
\emph{JC:} ``One estimate throughout (0.829 vs 0.830, 0.696 vs 0.697);
no raw and DenseCRF values mixed \ldots; \guillemotleft{}Eq. ??\guillemotright{} in Fig. 2; citations
for Attn2Mask, Grounded Diffusion and DatasetDM; every pointer to the
supplementary landing on the right section.''
\end{quote}

162 numbers traced: 137 match a committed result, 22 differ, 3 untraced.
Full list: \texttt{comparisons/notes/entry\_6\_consistency.md}.

{\def\LTcaptype{none} % do not increment counter
\begin{longtable}[]{@{}
  >{\raggedright\arraybackslash}p{(\linewidth - 4\tabcolsep) * \real{0.3333}}
  >{\raggedright\arraybackslash}p{(\linewidth - 4\tabcolsep) * \real{0.3333}}
  >{\raggedright\arraybackslash}p{(\linewidth - 4\tabcolsep) * \real{0.3333}}@{}}
\toprule\noalign{}
\begin{minipage}[b]{\linewidth}\raggedright
JC's item
\end{minipage} & \begin{minipage}[b]{\linewidth}\raggedright
Cause
\end{minipage} & \begin{minipage}[b]{\linewidth}\raggedright
Fix (location)
\end{minipage} \\
\midrule\noalign{}
\endhead
\bottomrule\noalign{}
\endlastfoot
0.829 vs 0.830, 0.696 vs 0.697 & three generations of the same pipeline
(v2, v3, a third) & use \textbf{v3: 0.829 / 0.696} everywhere; change
Sec. 5.3 (p.~6, \texttt{sec/5\_analysis.tex:35}) and Table 3 (p.~7,
\texttt{tables/tab\_hierarchy.tex:20}); Table 4's .830 is in-sample,
label it \\
raw and DenseCRF mixed & Sec. 5.2 (p.~5--6,
\texttt{sec/5\_analysis.tex:24}): COCO pair after DenseCRF, VOC pair
raw; abstract ``five points'' is raw while the rest is after DenseCRF;
Table 5 mixes raw and DenseCRF rows & one stage per comparison: ``0.818
vs 0.823, both after DenseCRF''; ``costs up to three points after
DenseCRF (five before)''; new Table 5 uses DenseCRF for all (\#3) \\
``Eq. ??'' in Fig. 2 (p.~5) & in JC's PDF the seed equation is split
into Eqs. (3) and (4); the caption's
\texttt{\textbackslash{}ref\{eq:weights\}} has no matching label in that
build & on Overleaf, put \texttt{\textbackslash{}label\{eq:weights\}} on
the seed equation (Eq. 4 in JC's PDF) and run LaTeX twice; check with a
search for ``??'' before sending \\
missing citations & no \texttt{references.bib} entries & Attn2Mask
(Yoshihashi et al., arXiv 2309.01369), Grounded Diffusion (Li et al.,
ICCV 2023), DatasetDM (Wu et al., NeurIPS 2023); BibTeX in
\texttt{comparisons/notes/entry\_venue\_and\_published\_block.md} \\
supplementary pointers & 10 pointers say only ``(supplementary
material)''; the supplement numbers its own Table 1, Fig. 2, Sec. 5,
which collide with the main paper's & point each to its section/table
(list in the note); number the supplement S1, S2, \ldots{} \\
\end{longtable}
}

\textbf{Slides:} \texttt{c6\_N.tex}.

\begin{center}\rule{0.5\linewidth}{0.5pt}\end{center}

\section{\#7 Table 1 --- deltas}

\begin{quote}
\emph{JC:} ``a COCO-Cap \ensuremath{\Delta}, the folds won (x/5) for each row \ldots{} and
make the reference of each \ensuremath{\Delta} explicit''; ``no COCO-Cap value at all in
the first three rows''; row 3 ``hides the largest jump (.691 \ensuremath{\rightarrow} .789)''.
\end{quote}

\subsection{\texorpdfstring{7.1 What is wrong with the current Table 1
(p.~4,
\texttt{tables/tab\_progression.tex})}{7.1 What is wrong with the current Table 1 (p.~4, tables/tab\_progression.tex)}}

\begin{itemize}
\tightlist
\item
  No COCO-Cap \ensuremath{\Delta}; rows 1--3 have no COCO-Cap value.
\item
  \ensuremath{\Delta} is not the difference between consecutive rows (each comes from its
  own paired run), e.g.~.789 \ensuremath{\rightarrow} .800 is +.011, printed +.012.
\item
  Row 3's \ensuremath{\Delta} (+.033) is measured against a row that is not in the table
  (.756, one pass through all 16 SA modules), which hides the .691 \ensuremath{\rightarrow}
  .789 jump.
\end{itemize}

\subsection{7.2 Proposed Table 1 (draft; the COCO-Cap values of rows
1--3 come from a fresh capture {[}pending{]})}

{\def\LTcaptype{none} % do not increment counter
\begin{longtable}[]{@{}
  >{\raggedright\arraybackslash}p{(\linewidth - 12\tabcolsep) * \real{0.1429}}
  >{\raggedright\arraybackslash}p{(\linewidth - 12\tabcolsep) * \real{0.1429}}
  >{\raggedright\arraybackslash}p{(\linewidth - 12\tabcolsep) * \real{0.1429}}
  >{\raggedright\arraybackslash}p{(\linewidth - 12\tabcolsep) * \real{0.1429}}
  >{\raggedright\arraybackslash}p{(\linewidth - 12\tabcolsep) * \real{0.1429}}
  >{\raggedright\arraybackslash}p{(\linewidth - 12\tabcolsep) * \real{0.1429}}
  >{\raggedright\arraybackslash}p{(\linewidth - 12\tabcolsep) * \real{0.1429}}@{}}
\toprule\noalign{}
\begin{minipage}[b]{\linewidth}\raggedright
Row
\end{minipage} & \begin{minipage}[b]{\linewidth}\raggedright
Change
\end{minipage} & \begin{minipage}[b]{\linewidth}\raggedright
VOC-Sim
\end{minipage} & \begin{minipage}[b]{\linewidth}\raggedright
\ensuremath{\Delta} VOC (vs row)
\end{minipage} & \begin{minipage}[b]{\linewidth}\raggedright
folds won
\end{minipage} & \begin{minipage}[b]{\linewidth}\raggedright
COCO-Cap
\end{minipage} & \begin{minipage}[b]{\linewidth}\raggedright
\ensuremath{\Delta} COCO (vs row)
\end{minipage} \\
\midrule\noalign{}
\endhead
\bottomrule\noalign{}
\endlastfoot
1 & CA only, all scales & .665 & & & {[}pending{]} & \\
2 & drop 32-px CA & .691 & +.026 (vs 1) & 4/5 & {[}pending{]} &
{[}pending{]} \\
3a & SA propagation, all 16 modules, one pass & .756 & +.065 (vs 2) &
{[}pending{]} & {[}pending{]} & {[}pending{]} \\
3b & SA 64 px only & .789 & +.033 (vs 3a) & 5/5 & {[}pending{]} &
{[}pending{]} \\
4 & two passes & .800 & +.011 (vs 3b) & 5/5\textsuperscript{b} & .625 & {[}pending{]} \\
5 & DenseCRF & .813 & +.013 (vs 4) & 5/5 & .637 & +.012 \\
6 & SA calls 10--20 & .817 & +.004 (vs 5) & 4/5 & .644 & +.007 \\
7 & SA-guided weights & .823 & +.006 (vs 6) & 4/5 & .662 &
\textbf{+.018} \\
8 & coarse-to-fine + cap & .829 & +.006 (vs 7) & 5/5 & .696 &
\textbf{+.034} \\
\end{longtable}
}

\textsuperscript{b} with DenseCRF. Folds won from supp. Tab. 10. \ensuremath{\Delta} printed as the
difference of the printed values; the paired \ensuremath{\Delta} and interval go in a
separate column or the supplement. The COCO-Cap \ensuremath{\Delta} of rows 7--8 needs the
threshold note of \#1 (+0.034 is two thresholds; +0.012 at a shared
one). LaTeX draft: \texttt{comparisons/drafts/table1\_jc7.tex}
{[}pending{]}. Putting it into \texttt{paper/} is F's.

\textbf{Slides:} F's \texttt{c7\_F.tex}.

\begin{center}\rule{0.5\linewidth}{0.5pt}\end{center}

\section{\#8 Paper structure and cost}

\begin{quote}
\emph{JC:} ``Rename Sec. 5 to \guillemotleft{}Experimental results\guillemotright{} and move Sec. 6
\ldots{} into it \ldots; Sec. 4 should include the hardware and software
details \ldots{} and the cost of the method: runtime and memory per
image.''
\end{quote}

\subsection{8.1 Structure}

{\def\LTcaptype{none} % do not increment counter
\begin{longtable}[]{@{}
  >{\raggedright\arraybackslash}p{(\linewidth - 2\tabcolsep) * \real{0.5000}}
  >{\raggedright\arraybackslash}p{(\linewidth - 2\tabcolsep) * \real{0.5000}}@{}}
\toprule\noalign{}
\begin{minipage}[b]{\linewidth}\raggedright
Now
\end{minipage} & \begin{minipage}[b]{\linewidth}\raggedright
Proposed
\end{minipage} \\
\midrule\noalign{}
\endhead
\bottomrule\noalign{}
\endlastfoot
4 Experimental setup & 4 Experimental setup: benchmarks, protocol (with
the override, \#1), metrics, \textbf{hardware, software, runtime and
memory} \\
5 What each stage contributes & \textbf{5 Experimental results} \\
6 Comparison with published methods & \textbf{5.1 Comparison} (new Table
5, \#3) --- first, so the narrowed claim is the first result \\
& 5.2--5.6 the former 5.1--5.5 (seed, envelope, cap, DenseCRF,
calibration) \\
& 5.7 Optimized tokens (\#4); 5.8 Other backbones (\#5) \\
7 Discussion and limitations & 6 Discussion and limitations (+ token
scope, CLIP-seeded methods, SDXL, SD3/FLUX out of scope) \\
Conclusion & 7 Conclusion \\
\end{longtable}
}

\subsection{8.2 Cost per image (measured; workplan entry to come)}

Configuration: 100 benchmark images (50 + 50), SD 1.5, 30 steps,
guidance 7.5, F's frozen runtime; median per image.

{\def\LTcaptype{none} % do not increment counter
\begin{longtable}[]{@{}
  >{\raggedright\arraybackslash}p{(\linewidth - 4\tabcolsep) * \real{0.3333}}
  >{\raggedright\arraybackslash}p{(\linewidth - 4\tabcolsep) * \real{0.3333}}
  >{\raggedright\arraybackslash}p{(\linewidth - 4\tabcolsep) * \real{0.3333}}@{}}
\toprule\noalign{}
\begin{minipage}[b]{\linewidth}\raggedright
Stage
\end{minipage} & \begin{minipage}[b]{\linewidth}\raggedright
A100 (40 GB)
\end{minipage} & \begin{minipage}[b]{\linewidth}\raggedright
V100
\end{minipage} \\
\midrule\noalign{}
\endhead
\bottomrule\noalign{}
\endlastfoot
generation, no attention recording & 2.66 s & {[}pending{]} \\
generation with attention recording & 6.94 s (+4.27 s) &
{[}pending{]} \\
seed + maps (hierarchy 0.04 s; same as fine-only) & 0.15 s &
{[}pending{]} \\
DenseCRF (CPU, one core) & 5.9 s & {[}pending{]} \\
\textbf{total, BRACE} (fine-only is the same) & \textbf{12.9 s} &
{[}pending{]} (earlier full run: median 17.3 s) \\
peak GPU memory & 5.1 GiB & {[}pending{]} (earlier: 5.3 GiB) \\
\end{longtable}
}

One-pass setting (\#5), per image on A100: SD 1.5 0.21 s / 5.1 GiB, SD
2.1 0.19 s / 5.5 GiB, SDXL 1.49 s / 15.6 GiB.

\textbf{Paragraph for Sec. 4 (draft):}

\begin{quote}
\emph{Hardware and cost.} Experiments ran on NVIDIA V100 (16/32 GB) and
A100 (40 GB) GPUs with PyTorch 2.4.1 (CUDA 12.1), diffusers 0.36 and
transformers 5.6. Per image on an A100, generation alone takes 2.7 s;
recording the attention maps adds 4.3 s; building the seed and the
hierarchy takes 0.15 s, the same as fine-only propagation; DenseCRF
takes 5.9 s on one CPU core. Peak GPU memory is 5.1 GiB. The hierarchy
therefore costs almost nothing over fine-only propagation; recording the
attention and DenseCRF dominate.
\end{quote}

(The paragraph says ``recording'' rather than ``storing 64\ensuremath{\times}64 SA maps
over many calls'', because the maps are reduced on the fly; JC's concern
about memory is answered by the 5.1 GiB peak.)

\textbf{Slides:} \texttt{c8\_N.tex} (structure); cost frame to add.

\begin{center}\rule{0.5\linewidth}{0.5pt}\end{center}

\section{\#9 Target venue: ICPR 2027 --- done}

{\def\LTcaptype{none} % do not increment counter
\begin{longtable}[]{@{}
  >{\raggedright\arraybackslash}p{(\linewidth - 2\tabcolsep) * \real{0.5000}}
  >{\raggedright\arraybackslash}p{(\linewidth - 2\tabcolsep) * \real{0.5000}}@{}}
\toprule\noalign{}
\begin{minipage}[b]{\linewidth}\raggedright
Item
\end{minipage} & \begin{minipage}[b]{\linewidth}\raggedright
ICPR 2027 (ICPRv)
\end{minipage} \\
\midrule\noalign{}
\endhead
\bottomrule\noalign{}
\endlastfoot
Registration (author list fixed) & \textbf{1 March 2027}, AoE \\
Submission & \textbf{8 March 2027}, AoE \\
Notification / camera-ready & 14 June / 28 June 2027 \\
Format & Springer LNCS, one column, \textbf{15 pages + references} \\
Review & double-blind; mandatory self-review pages incl.~an LLM-use
disclosure \\
Conference & online only, 4--6 and 14--15 Oct 2027 \\
Preprints & arXiv allowed \\
\end{longtable}
}

\textbf{Page fit.} The current text, converted to the ICPR template
unchanged, fills about \textbf{13.5 of 15 pages} before any of JC's
additions (new Table 5, backbones, tokens, cost, Table 1 columns).
Needed: the \texttt{llncs} class, table captions above tables, Table 3
restructured (6.5 cm too wide). Move to the supplement: Table 4 (one
sentence instead), Fig. 5, one of Figs. 3 and 6, the top half of Table
2.

\textbf{To decide with JC:} author list and order (fixed on 1 March);
online-only venue OK; long paper (15 pages).

\textbf{Slides:} \texttt{c9\_N.tex}.

\begin{center}\rule{0.5\linewidth}{0.5pt}\end{center}

\section{Summary for the meeting}

{\def\LTcaptype{none} % do not increment counter
\begin{longtable}[]{@{}
  >{\raggedright\arraybackslash}p{(\linewidth - 4\tabcolsep) * \real{0.3333}}
  >{\raggedright\arraybackslash}p{(\linewidth - 4\tabcolsep) * \real{0.3333}}
  >{\raggedright\arraybackslash}p{(\linewidth - 4\tabcolsep) * \real{0.3333}}@{}}
\toprule\noalign{}
\begin{minipage}[b]{\linewidth}\raggedright
\#
\end{minipage} & \begin{minipage}[b]{\linewidth}\raggedright
Status
\end{minipage} & \begin{minipage}[b]{\linewidth}\raggedright
One-line answer
\end{minipage} \\
\midrule\noalign{}
\endhead
\bottomrule\noalign{}
\endlastfoot
1 & text drafted & Reframe around clutter and calibration; disclose the
override and its timing; drop ``six times''; narrow to diffusion-only \\
2 & done & Contributions 1--2 well supported; 3--4 rely on the
supplement; add the rule's choice next to ours \\
3 & done (OVAM, DAAM, SCLIP {[}pending{]}) & BRACE best among
diffusion-only methods; NERVE (with CLIP) better on COCO-Cap \\
4 & done (corrected-prompt rerun {[}pending{]}) & No test leakage; the
token helps VOC-Sim, and with it the hierarchy no longer helps on
COCO-Cap \\
5 & done & Hierarchy holds on SD 1.5 and SD 2.1, fails on SDXL; GoCA
wins on SDXL \\
6 & done & 22 fixes listed \\
7 & numbers {[}pending{]} & New columns: COCO \ensuremath{\Delta}, folds won, explicit
references, visible .691 \ensuremath{\rightarrow} .756 \ensuremath{\rightarrow} .789 \\
8 & done (V100 {[}pending{]}) & Structure proposal; 12.9 s and 5.1 GiB
per image; the hierarchy is nearly free \\
9 & done & ICPR confirmed; about 13.5 of 15 pages used before the
additions \\
\end{longtable}
}

\end{document}
